% Please do not change %
\documentclass[11pt]{article} % font size
\usepackage[margin=1in]{geometry} % margins
\usepackage{amsmath,amsthm,amssymb} % for the  common "math symbols"
\usepackage{algorithm}
\usepackage{algpseudocode}
\usepackage[shortlabels]{enumitem} % for enumeration
\usepackage{booktabs} % if you need tables
\usepackage{listings}
\usepackage{color}
\usepackage{quiver}
\usepackage{ulem}
\usepackage{bbm}
\usepackage{hyperref}
\DeclareMathOperator*{\argmax}{argmax}
\DeclareMathOperator*{\argmin}{argmin}

% \theoremstyle{definition}
\newtheorem{theorem}{Theorem}[section]
\newtheorem{corollary}[theorem]{Corollary}
\newtheorem{lemma}[theorem]{Lemma}
\newtheorem{proposition}[theorem]{Proposition}
\newtheorem{claim}[theorem]{Claim}
\newtheorem{conjecture}[theorem]{Conjecture}
\theoremstyle{definition}
\newtheorem{definition/}[theorem]{Definition}
\theoremstyle{definition}
\newtheorem{fact}[theorem]{Fact}
\theoremstyle{definition}
\newtheorem{example/}[theorem]{Example}
\newtheorem{note}[theorem]{Note}

\newenvironment{example}{
  \renewcommand{\qedsymbol}{$\blacktriangle$}%
  \pushQED{\qed}\begin{example/}
}{\popQED\end{example/}}

\newenvironment{definition}{
  \renewcommand{\qedsymbol}{$\blacktriangle$}%
  \pushQED{\qed}\begin{definition/}}
{\popQED\end{definition/}}

\renewcommand{\thesubsection}{\alph{subsection}}

% Feel free to include additional packages (if needed), but make sure it does not override the margin/font requirements.
%
\setlength{\parindent}{0pt} % no indent for new paragraphs.
\setlength{\parskip}{5pt} % distance between paragraphs, which will be used when you have a blank line in your LaTeX code.


%%%%%%%%%%%%%%%%%%%%%%%%%%%%%%%%%%%%%%%%%%%%%%%%%%%%%%%%%%%%%%%%%%%%%%%%%%%%%%
\title{Categories, Proofs, and Processes: \\Mini-Project}
%
\author{1081502}
\date{Hilary Term}
%
\begin{document}
%
\maketitle


%\begin{table}[h]
%\caption{A simple table.}
%\centering
%\begin{tabular}[t]{lrrrr}
%\toprule
%Graphs & Can & Be  & Fun & Too \\ 
%\midrule 
%$G_1$ & $1$  & $2$ & $3$ & $4$ \\ 
%$G_2$ & $-1$  & $-2$ & $-3$ & $-4$ \\ 
%$G_3$ & $0.1$  & $0.2$ & $0.3$ & $0.4$ \\ 
%\bottomrule
%\end{tabular}
%\label{tab}
%\end{table}

% A functor $F:\mathcal{C} \rightarrow \mathcal{D}$ is
% \begin{itemize}
%     \item \textbf{faithful} if each map $F_{A,B}:\mathcal{C}(A,B) \rightarrow \mathcal{D}(F\ A, F\ B)$ is injective;
%     \item \textbf{full} if each map $F_{A,B}:\mathcal{C}(A,B) \rightarrow \mathcal{D}(F\ A, F\ B)$ is surjective
% \end{itemize}

\section{}
\subsection{}
\begin{definition}[\cite{Abramsky_2010}]
A \underline{natural isomorphism} is a natural transformation $t : F \rightarrow G$ 
where each arrow $t_A$ in its family of arrows is isomorphic:
\[t_A : F\ A \xrightarrow[]{\cong} G\ A\] 
\end{definition}

\begin{definition}[\cite{maclane:71}, \cite{nlab:bicategory} \cite{BénabouJean2006Itb}]
A \textbf{\underline{bicategory}} $B$ consists of
\begin{itemize}
    \item A collection of objects $a,b,c,...$ called \textbf{\textit{0-cells}}.
    \item For each pair of 0-cells, $a,b$, a category $B(a,b)$ where its
    \begin{itemize}
        \item Objects are morphisms between $a,b$ called \textbf{\textit{1-cells}}.
        \item Morphisms between objects are called \textbf{\textit{2-cells}}, or equivalently, natural transformations.
    \end{itemize}
    % https://q.uiver.app/#q=WzAsMixbMCwwLCJhIl0sWzEsMCwiYiJdLFswLDEsImYiLDAseyJvZmZzZXQiOi01fV0sWzAsMSwiZyIsMix7Im9mZnNldCI6NX1dLFsyLDMsIlxcZXRhIiwwLHsic2hvcnRlbiI6eyJzb3VyY2UiOjIwLCJ0YXJnZXQiOjIwfX1dXQ==
\[\begin{tikzcd}
	a & b
	\arrow[""{name=0, anchor=center, inner sep=0}, "f", shift left=5, from=1-1, to=1-2]
	\arrow[""{name=1, anchor=center, inner sep=0}, "g"', shift right=5, from=1-1, to=1-2]
	\arrow["\eta", shorten <=3pt, shorten >=3pt, Rightarrow, from=0, to=1]
\end{tikzcd}\]
    \item For each 0-cell $a$, the 1-cell $1_a \in B(a,a)$ called the identity cell at $a$.
    \item For each triple of 0-cells, $a,b,c$, a bifunctor 
    \[\circ: B(b,c) \times B(a,b) \rightarrow B(a,c)\]
    called \textbf{\textit{horizontal composition}}.
    \item For each 1-cell $f:a \rightarrow b$, the 2-cell identity $1_f:f \Rightarrow f$ called the 2-identity.
    \item For each triple of 1-cells, $f,g,h:a \rightarrow b$
    with 2-cells $\eta: f\Rightarrow g, \theta: g\Rightarrow h$, a vertical composite 
    \[\theta \bullet \eta : f \Rightarrow h\]
% https://q.uiver.app/#q=WzAsNCxbMCwwLCJhIl0sWzEsMCwiYiJdLFszLDAsImEiXSxbNCwwLCJiIl0sWzAsMSwiZiIsMCx7Im9mZnNldCI6LTV9XSxbMCwxLCJoIiwyLHsib2Zmc2V0Ijo1fV0sWzAsMSwiZyIsMV0sWzIsMywiZiIsMCx7Im9mZnNldCI6LTV9XSxbMiwzLCJoIiwyLHsib2Zmc2V0Ijo1fV0sWzYsNSwiXFx0aGV0YSIsMCx7InNob3J0ZW4iOnsic291cmNlIjoyMCwidGFyZ2V0IjoyMH19XSxbNCw2LCJcXGV0YSIsMCx7InNob3J0ZW4iOnsic291cmNlIjoyMCwidGFyZ2V0IjoyMH19XSxbNyw4LCJcXHRoZXRhXFxidWxsZXQgXFxldGEiLDAseyJzaG9ydGVuIjp7InNvdXJjZSI6MjAsInRhcmdldCI6MjB9fV1d
\[\begin{tikzcd}
	a & b && a & b
	\arrow[""{name=0, anchor=center, inner sep=0}, "f", shift left=5, from=1-1, to=1-2]
	\arrow[""{name=1, anchor=center, inner sep=0}, "h"', shift right=5, from=1-1, to=1-2]
	\arrow[""{name=2, anchor=center, inner sep=0}, "g"{description}, from=1-1, to=1-2]
	\arrow[""{name=3, anchor=center, inner sep=0}, "f", shift left=5, from=1-4, to=1-5]
	\arrow[""{name=4, anchor=center, inner sep=0}, "h"', shift right=5, from=1-4, to=1-5]
	\arrow["\theta", shorten <=1pt, shorten >=1pt, Rightarrow, from=2, to=1]
	\arrow["\eta", shorten <=1pt, shorten >=1pt, Rightarrow, from=0, to=2]
	\arrow["{\theta\bullet \eta}", shorten <=3pt, shorten >=3pt, Rightarrow, from=3, to=4]
\end{tikzcd}\]
    \item For each triple of 0-cells $a,b,c$, 1-cells $f,g:a\rightarrow b,\quad h:b\rightarrow c$, 
    2-cell $\eta : f \Rightarrow g$, a \textbf{\textit{left whiskering}}
    \[h \triangleleft \eta: h\circ f \Rightarrow h\circ g\]
    \item For each triple of 0-cells $a,b,c$, 1-cells $f:a\rightarrow b,\quad g,h:b\rightarrow c$, 
    2-cell $\eta : g \Rightarrow h$, a \textbf{\textit{right whiskering}}
    \[\eta \triangleright f: g\circ f \Rightarrow h\circ f\]
    \item Natural isomorphisms $\rho,\lambda$ where for each 1-cell $f: a\rightarrow b$, 
    components $\rho_f,\lambda_f$ called \textbf{\textit{right and left unitors}} respectively where
    \[\rho_f : f \circ 1_a \Rightarrow f, \quad \lambda_f : 1_b \circ f \Rightarrow f\]
    \item A natural isomorphism $\alpha$, 
    where for each triple of 1-cells $a\xrightarrow[]{f} b\xrightarrow[]{g} c\xrightarrow[]{h} d$,
    the component $\alpha_{h,g,f}$ called an \textbf{\textit{associator}} defined as
    \[\alpha_{h,g,f} : h\circ(g\circ f) \Rightarrow (h \circ g) \circ f\]
\end{itemize}
such that the following conditions hold
\begin{itemize}
    \item The \textit{pentagon identity} is satisfied by the associators - for any quadruple of 1-cells
    \[a\xrightarrow[]{f} b\xrightarrow[]{g} c\xrightarrow[]{h} d \xrightarrow[]{k} e\]
    the following pentagon must commute.
    % https://q.uiver.app/#q=WzAsNSxbMCwwLCJrXFxjaXJjKGhcXGNpcmMoZ1xcY2lyYyBmKSkpIl0sWzEsMCwiKGtcXGNpcmMgaCkgXFxjaXJjIChnXFxjaXJjIGYpIl0sWzAsMSwia1xcY2lyYyAoKGhcXGNpcmMgZylcXGNpcmMgZikiXSxbMiwwLCIoKGtcXGNpcmMgaCkgXFxjaXJjIGcpIFxcY2lyYyBmKSJdLFsyLDEsIihrXFxjaXJjIChoXFxjaXJjIGcpKSBcXGNpcmMgZiJdLFswLDEsIlxcYWxwaGFfe2ssIGgsIGdcXGNpcmMgZn0iXSxbMCwyLCIxX2sgXFx0cmlhbmdsZWxlZnQgXFxhbHBoYV97aCxnLGZ9IiwyXSxbMSwzLCJcXGFscGhhX3trXFxjaXJjIGgsIGcsIGZ9Il0sWzMsNCwiXFxhbHBoYV97ayxoLGd9IFxcdHJpYW5nbGVyaWdodCAxX2YiXSxbMiw0LCJcXGFscGhhX3trLGhcXGNpcmMgZywgZn0iLDJdXQ==
\[\begin{tikzcd}
	{k\circ(h\circ(g\circ f)))} & {(k\circ h) \circ (g\circ f)} & {((k\circ h) \circ g) \circ f)} \\
	{k\circ ((h\circ g)\circ f)} && {(k\circ (h\circ g)) \circ f}
	\arrow["{\alpha_{k, h, g\circ f}}", from=1-1, to=1-2]
	\arrow["{1_k \triangleleft \alpha_{h,g,f}}"', from=1-1, to=2-1]
	\arrow["{\alpha_{k\circ h, g, f}}", from=1-2, to=1-3]
	\arrow["{\alpha_{k,h,g} \triangleright 1_f}", from=1-3, to=2-3]
	\arrow["{\alpha_{k,h\circ g, f}}"', from=2-1, to=2-3]
\end{tikzcd}\]
    \item The \textit{triangle identity} is satisfied by the unitors - 
    for any pair of 1-cells $a\xrightarrow[]{f} b\xrightarrow[]{g} c$, the following triangle must commute.
    % https://q.uiver.app/#q=WzAsMyxbMCwwLCJnXFxjaXJjKDFfYlxcY2lyYyBmKSJdLFswLDEsIihnXFxjaXJjIDFfYikgXFxjaXJjIGYiXSxbMSwxLCJnXFxjaXJjIGYiXSxbMCwxLCJcXGFscGhhX3tnLDFfYixmfSIsMl0sWzAsMiwiMV9nXFx0cmlhbmdsZWxlZnRcXGxhbWJkYV97YSxifSJdLFsxLDIsIlxccmhvX3thLGJ9XFx0cmlhbmdsZXJpZ2h0IDFfZiIsMl1d
\[\begin{tikzcd}
	{g\circ(1_b\circ f)} \\
	{(g\circ 1_b) \circ f} & {g\circ f}
	\arrow["{\alpha_{g,1_b,f}}"', from=1-1, to=2-1]
	\arrow["{1_g\triangleleft\lambda_{a,b}}", from=1-1, to=2-2]
	\arrow["{\rho_{a,b}\triangleright 1_f}"', from=2-1, to=2-2]
\end{tikzcd}\]
\end{itemize}
A bicategory is also called a \textit{weak 2-category}.
\end{definition}

\clearpage
\subsection{}
\begin{definition}[Monoidal Category \cite{Abramsky_2010}]
\label{def:moncat}
  A \textbf{\underline{monoidal category}} is a structure $(\mathcal{C}, \otimes, I, \alpha, \lambda, \rho)$ where
\begin{itemize}
  \item $\mathcal{C}$ is a category,
  \item $\otimes: \mathcal{C} \times \mathcal{C} \to \mathcal{C}$ is a functor (\textit{monoidal product}),
  \item $I$ is a distinguished object of $\mathcal{C}$ (\textit{unit})
  \item $\alpha, \lambda, \rho$ are natural isomorphisms (\textit{associator, left unitor, right unitor}) 
  with components
  \[\alpha_{A,B,C} : A\otimes (B \otimes C) \xrightarrow[]{\cong} (A\otimes B) \otimes C\]
  \[\lambda_A : I \otimes A \xrightarrow[]{\cong} A, \quad \rho_A : A\otimes I \xrightarrow[]{\cong} A\]
  such that $\lambda_I = \rho_I : I\otimes I \to I$ and the following diagrams commute.
% https://q.uiver.app/#q=WzAsOCxbMCwwLCJBXFxvdGltZXMgKEJcXG90aW1lcyhDXFxvdGltZXMgRCkpKSJdLFsxLDAsIihBXFxvdGltZXMgQilcXG90aW1lcyAoQ1xcb3RpbWVzIEQpIl0sWzIsMCwiKChBXFxvdGltZXMgQilcXG90aW1lcyBDKVxcb3RpbWVzIEQpIl0sWzAsMSwiQVxcb3RpbWVzICgoQlxcb3RpbWVzIEMpIFxcb3RpbWVzIEQpIl0sWzIsMSwiKEFcXG90aW1lcyAoQlxcb3RpbWVzIEMpKSBcXG90aW1lcyBEIl0sWzAsMiwiQVxcb3RpbWVzIChJXFxvdGltZXMgQikiXSxbMCwzLCIoQVxcb3RpbWVzIEkpIFxcb3RpbWVzIEIiXSxbMSwzLCJBXFxvdGltZXMgQiJdLFswLDEsIlxcYWxwaGFfe0EsQixDXFxvdGltZXMgRH0iXSxbMSwyLCJcXGFscGhhX3tBXFxvdGltZXMgQiwgQyxEfSJdLFswLDMsIjFfQSBcXG90aW1lcyBcXGFscGhhX3tCLEMsRH0iLDJdLFszLDQsIlxcYWxwaGFfe0EsQlxcb3RpbWVzIEMsRH0iLDJdLFsyLDQsIlxcYWxwaGFfe0EsQixDfVxcb3RpbWVzIDFfRCJdLFs1LDYsIlxcYWxwaGFfe0EsSSxCfSJdLFs1LDcsIjFfQSBcXG90aW1lcyBcXGxhbWJkYV9CIl0sWzYsNywiXFxyaG9fQVxcb3RpbWVzIDFfQiIsMl1d
\[\begin{tikzcd}
	{A\otimes (B\otimes(C\otimes D)))} & {(A\otimes B)\otimes (C\otimes D)} & {((A\otimes B)\otimes C)\otimes D)} \\
	{A\otimes ((B\otimes C) \otimes D)} && {(A\otimes (B\otimes C)) \otimes D} \\
	{A\otimes (I\otimes B)} \\
	{(A\otimes I) \otimes B} & {A\otimes B}
	\arrow["{\alpha_{A,B,C\otimes D}}", from=1-1, to=1-2]
	\arrow["{\alpha_{A\otimes B, C,D}}", from=1-2, to=1-3]
	\arrow["{1_A \otimes \alpha_{B,C,D}}"', from=1-1, to=2-1]
	\arrow["{\alpha_{A,B\otimes C,D}}"', from=2-1, to=2-3]
	\arrow["{\alpha_{A,B,C}\otimes 1_D}", from=1-3, to=2-3]
	\arrow["{\alpha_{A,I,B}}", from=3-1, to=4-1]
	\arrow["{1_A \otimes \lambda_B}", from=3-1, to=4-2]
	\arrow["{\rho_A\otimes 1_B}"', from=4-1, to=4-2]
\end{tikzcd}\]
\end{itemize}
\end{definition}
\begin{definition}
  \label{def:strict-mon}
  A monoidal category is \underline{strict} if $\alpha,\lambda,\rho$ are all identities.
\end{definition}
\begin{definition}[\cite{maclane:71}, \cite{nlab:bicategory}]
A monoidal category $M$ is a bicategory $BM$ with a single object $\bullet$ where
\begin{itemize}
    \item Objects $A$ of $M$ become 1-cells $[A]:\bullet \rightarrow \bullet$ of $BM$
    \item Composition of 1-cells in $BM$ follow the definition
    \[[A] \circ [B] := [A\otimes B]\]
    using the monoidal product $\otimes$ of $M$.
    \item The monoidal unit $I$ in $M$ becomes 
    the identity cell $1_{\bullet} := [I] : \bullet \rightarrow \bullet$ in $BM$.
    \item Morphisms $f:A\rightarrow B$ in $M$ become 2-cells $[f] : [A] \Rightarrow [B]$ in $BM$.
    \item Identity morphisms $1_A$ in $M$ become 2-identities $[1_A]$ in $BM$.
    \item Natural isomorphism $\alpha$ of $M$ where 
    \[\alpha_{A,B,C} : (A\otimes B) \otimes C \Rightarrow A\otimes (B\otimes C)\]
    becomes the associator 
    \[\alpha_{[A],[B],[C]} : ([A] \circ [B]) \circ [C] \Rightarrow [A] \circ ([B] \circ [C])\]
    for the natural isomorphism $\alpha$ of $BM$
    \item Natural isomorphisms $\lambda, \rho$ where each component
    \[\lambda_A : I \otimes A \Rightarrow A, \quad \rho_A: A \otimes I \Rightarrow A\]
    of $M$ respectively becomes the left and right unitors
    \[\lambda_{[A]} : [I] \circ [A] \Rightarrow [A], \quad
    \rho_{[A]} : [A] \circ [I] \Rightarrow [A]\]
    for the natural isomorphisms $\lambda, \rho$ of $BM$.
\end{itemize}

\end{definition}

The coherence conditions then hold, where the following diagrams commute in $BM$ 
iff the corresponding diagrams \cite{Abramsky_2010} commute in $M$.
\begin{itemize}
\item Unitors
% https://q.uiver.app/#q=WzAsNyxbMCwwLCJbQV1cXGNpcmMoW0ldXFxjaXJjIFtCXSkiXSxbMCwxLCIoW0FdIFxcY2lyYyBbSV0pIFxcY2lyYyBbQl0iXSxbMSwxLCJbQV1cXGNpcmMgW0JdIl0sWzIsMSwiXFxpZmYiXSxbMywwLCJBXFxvdGltZXMgKElcXG90aW1lcyBCKSJdLFszLDEsIihBXFxvdGltZXMgSSkgXFxvdGltZXMgQiJdLFs0LDEsIkFcXG90aW1lcyBCIl0sWzAsMSwiXFxhbHBoYV97W0JdLFtJXSxbQV19IiwyXSxbMCwyLCIxX3tbQV19XFx0cmlhbmdsZWxlZnRcXGxhbWJkYV97W0JdfSJdLFsxLDIsIlxccmhvX3tbQV19XFx0cmlhbmdsZXJpZ2h0IDFfe1tCXX0iLDJdLFs0LDUsIlxcYWxwaGFfe0EsSSxCfSIsMl0sWzQsNiwiMV9BIFxcb3RpbWVzIFxcbGFtYmRhX0IiXSxbNSw2LCJcXHJob19BXFxvdGltZXMxX0IiLDJdXQ==
\[\begin{tikzcd}
	{[A]\circ([I]\circ [B])} &&& {A\otimes (I\otimes B)} \\
	{([A] \circ [I]) \circ [B]} & {[A]\circ [B]} & \iff & {(A\otimes I) \otimes B} & {A\otimes B}
	\arrow["{\alpha_{[B],[I],[A]}}"', from=1-1, to=2-1]
	\arrow["{1_{[A]}\triangleleft\lambda_{[B]}}", from=1-1, to=2-2]
	\arrow["{\rho_{[A]}\triangleright 1_{[B]}}"', from=2-1, to=2-2]
	\arrow["{\alpha_{A,I,B}}"', from=1-4, to=2-4]
	\arrow["{1_A \otimes \lambda_B}", from=1-4, to=2-5]
	\arrow["{\rho_A\otimes1_B}"', from=2-4, to=2-5]
\end{tikzcd}\]
\item Associators
% https://q.uiver.app/#q=WzAsMTEsWzAsMCwiW0FdXFxjaXJjKFtCXVxcY2lyYyhbQ11cXGNpcmMgW0RdKSkpIl0sWzEsMCwiKFtBXVxcY2lyYyBbQl0pIFxcY2lyYyAoW0NdXFxjaXJjIFtEXSkiXSxbMCwxLCJbQV1cXGNpcmMgKChbQl1cXGNpcmMgW0NdKVxcY2lyYyBbRF0pIl0sWzIsMCwiKChbQV1cXGNpcmMgW0JdKSBcXGNpcmMgW0NdKSBcXGNpcmMgW0RdKSJdLFsyLDEsIihbQV1cXGNpcmMgKFtCXVxcY2lyYyBbQ10pKSBcXGNpcmMgW0RdIl0sWzEsMiwiXFxpZmYiXSxbMCwzLCJBXFxvdGltZXMgKEJcXG90aW1lcyhDXFxvdGltZXMgRCkpKSJdLFsxLDMsIihBXFxvdGltZXMgQilcXG90aW1lcyAoQ1xcb3RpbWVzIEQpIl0sWzIsMywiKChBXFxvdGltZXMgQilcXG90aW1lcyBDKVxcb3RpbWVzIEQpIl0sWzAsNCwiQVxcb3RpbWVzICgoQlxcb3RpbWVzIEMpIFxcb3RpbWVzIEQpIl0sWzIsNCwiKEFcXG90aW1lcyAoQlxcb3RpbWVzIEMpKSBcXG90aW1lcyBEIl0sWzAsMSwiXFxhbHBoYV97W0FdLCBbQl0sIFtDXVxcY2lyYyBbRF19Il0sWzAsMiwiMV97W0FdfSBcXHRyaWFuZ2xlbGVmdCBcXGFscGhhX3tbQl0sW0NdLFtEXX0iLDJdLFsxLDMsIlxcYWxwaGFfe1tBXVxcY2lyYyBbQl0sW0NdLFtEXX0iXSxbMyw0LCJcXGFscGhhX3tbQV0sW0JdLFtDXX0gXFx0cmlhbmdsZXJpZ2h0IDFfe1tEXX0iXSxbMiw0LCJcXGFscGhhX3tbQV0sW0JdXFxjaXJjIFtDXSwgW0RdfSIsMl0sWzYsNywiXFxhbHBoYV97QSxCLENcXG90aW1lcyBEfSJdLFs3LDgsIlxcYWxwaGFfe0FcXG90aW1lcyBCLCBDLER9Il0sWzYsOSwiMV9BIFxcb3RpbWVzIFxcYWxwaGFfe0IsQyxEfSIsMl0sWzksMTAsIlxcYWxwaGFfe0EsQlxcb3RpbWVzIEMsRH0iLDJdLFs4LDEwLCJcXGFscGhhX3tBLEIsQ31cXG90aW1lcyAxX0QiXV0=
\[\begin{tikzcd}
	{[A]\circ([B]\circ([C]\circ [D])))} & {([A]\circ [B]) \circ ([C]\circ [D])} & {(([A]\circ [B]) \circ [C]) \circ [D])} \\
	{[A]\circ (([B]\circ [C])\circ [D])} && {([A]\circ ([B]\circ [C])) \circ [D]} \\
	& \iff \\
	{A\otimes (B\otimes(C\otimes D)))} & {(A\otimes B)\otimes (C\otimes D)} & {((A\otimes B)\otimes C)\otimes D)} \\
	{A\otimes ((B\otimes C) \otimes D)} && {(A\otimes (B\otimes C)) \otimes D}
	\arrow["{\alpha_{[A], [B], [C]\circ [D]}}", from=1-1, to=1-2]
	\arrow["{1_{[A]} \triangleleft \alpha_{[B],[C],[D]}}"', from=1-1, to=2-1]
	\arrow["{\alpha_{[A]\circ [B],[C],[D]}}", from=1-2, to=1-3]
	\arrow["{\alpha_{[A],[B],[C]} \triangleright 1_{[D]}}", from=1-3, to=2-3]
	\arrow["{\alpha_{[A],[B]\circ [C], [D]}}"', from=2-1, to=2-3]
	\arrow["{\alpha_{A,B,C\otimes D}}", from=4-1, to=4-2]
	\arrow["{\alpha_{A\otimes B, C,D}}", from=4-2, to=4-3]
	\arrow["{1_A \otimes \alpha_{B,C,D}}"', from=4-1, to=5-1]
	\arrow["{\alpha_{A,B\otimes C,D}}"', from=5-1, to=5-3]
	\arrow["{\alpha_{A,B,C}\otimes 1_D}", from=4-3, to=5-3]
\end{tikzcd}\]
\end{itemize}

\clearpage
\section{}

\subsection{… , a monoid is a category, …}
First, we must list the definitions of both a monoid and category.
\begin{definition}[Monoid \cite{Abramsky_2010}]
  \label{def:monoid}
    A \textbf{\underline{monoid}} is a structure $(M,\cdot, 1)$ where $M$ is a set,
    \[- \cdot - : M\times M \rightarrow M\]
    is a binary operation, and $1 \in M$, satisfying the axioms $\forall x,y,z\in M$.
    \[(x\cdot y)\cdot z = x\cdot (y\cdot z),\quad 1\cdot x = x = x\cdot 1\]
\end{definition}
\begin{definition}[Category \cite{Abramsky_2010} \cite{nlab:category} \cite{Grothendieck:61}]
A \textbf{\underline{category}} consists of
\begin{itemize}
    \item A collection Ob$(\mathcal{C})$ of \textbf{\textit{objects}}.
    \item A collection Ar$(\mathcal{C})$ of \textbf{\textit{arrows}} or morphisms.
    \item For every arrow $f$, an object $s(f)$ called its source or domain, 
    and an object $t(f)$ called its target or codomain. 
    An arrow with domain $s(f)=A$ and codomain $t(f)=B$ is denoted $f:A \rightarrow B$.
    \item For each pair of objects $A,B$, we define the set 
    \[\mathcal{C}(A,B) := \{f \in \text{Ar}(\mathcal{C}) | f: A\rightarrow B\}\]
    called the \textbf{\textit{hom-set}}.
    \item For each pair of morphisms $f$ in $\mathcal{C}(A,B)$, and $g$ in $\mathcal{C}(B,C)$, 
    a composite morphism  $g\circ f$ in $\mathcal{C}(A,C)$.
    \[A\xrightarrow[]{f} B \xrightarrow[]{g} C\]
    \item For each object $A$, an \textbf{\textit{identity}} arrow $1_A : A \rightarrow A$
\end{itemize}
where the following axioms are satisfied 
whenever the domains and codomains match so that compositions are well-defined.
\begin{itemize}
    \item Associativity
    \[\forall f,g,h\in\text{Ar}(\mathcal{C}), h\circ (g\circ f) = (h\circ g)\circ f\]
    \item Unit
    \[\forall f:A\rightarrow B, f\circ 1_A = f = 1_B \circ f\]
\end{itemize}
\end{definition}

\clearpage
A monoid $(M,\cdot,1_M)$ can then be defined as a category $\mathcal{M}$ with one object $*$ where
\begin{itemize}
    \item Each arrow $f_e:*\rightarrow *$ of the category corresponds to each element $e\in M$ of the monoid.
    \item The source and target of all arrows are the same object $*$.
    \item There is a single hom-set $\mathcal{M}(*, *)$.
    \item Composition $f_b \circ f_a$ corresponds to applying binary operations $-\cdot a \cdot b$
    \item The identity arrow on $*, 1_{*}$ corresponds to the unit of the monoid $1_M\in M$.
\end{itemize}
so the axioms are satisfied in the category $\mathcal{M}$ 
iff the axioms are satisfied in the monoid $(M,\cdot,1_M)$.
\begin{itemize}
    \item Associativity
    \[\forall f_a,f_b,f_c\in\text{Ar}(\mathcal{M}), f_c\circ (f_b\circ f_a) = (f_c\circ f_b)\circ f_a
    \iff \forall a,b,c\in M, (a \cdot b) \cdot c = a \cdot (b\cdot c)\]
    \item Unit
    \[\forall f_a\in\text{Ar}(\mathcal{M}), f_a\circ 1_{*} = f_a = 1_{*} \circ f_a
    \iff \forall a\in M, 1_M \cdot a = a = a \cdot 1_M\]
\end{itemize}

So every monoid is a special kind of category, and categories are the more general concept.

\clearpage
\subsection{… , a category is a monad, …}
We note the following definitions.
\begin{definition}[Strict]
  A \underline{strict} bicategory is a bicategory where the natural isomorphisms $\alpha,\rho,\lambda$ are identities.
\end{definition}
\begin{definition}[Strict 2-category \cite{Lack_2009}]
  \textbf{\underline{Strict 2-categories}} are strict bicategories. Conversely, bicategories are weak 2-categories.
\end{definition}
We can note that bicategories are then the generalisation of strict 2-categories - 
the notion of a \textit{general 2-category}.
\begin{note}
  A 2-category in this paper refers to any strict or general 2-category, including bicategories.
\end{note}

One way to define the concept of monads is the following.
\begin{note}
  \[T^2 := T \circ T, \quad T^3 := T \circ T^2\]
\end{note}
\begin{definition}[Monad \cite{Abramsky_2010} \cite{Sobociński:2017}]
\label{def:monad}
A monad over a category $\mathcal{C}$ is a triple $(T,\eta,\mu)$ 
where $T:\mathcal{C}\rightarrow\mathcal{C}$ is an endofunctor and 
$\eta : 1_{\mathcal{C}} \rightarrow T, \mu : T^2 \rightarrow T$ are natural transformations
such that the following diagram commutes.
% https://q.uiver.app/#q=WzAsOCxbMCwwLCJUXjMgQSJdLFsxLDAsIlReMkEiXSxbMCwxLCJUXjJBIl0sWzEsMSwiVEEiXSxbMywwLCJUQSJdLFs0LDAsIlReMkEiXSxbMywxLCJUXjIgQSJdLFs0LDEsIlRBIl0sWzAsMSwiXFxtdV97VEF9Il0sWzAsMiwiVFxcbXVfQSIsMl0sWzEsMywiXFxtdV9BIl0sWzIsMywiXFxtdV9BIiwyXSxbNCw1LCJcXGV0YV97VEF9Il0sWzQsNiwiVF97XFxldGEgQX0iLDJdLFs2LDcsIlxcbXVfQSIsMl0sWzUsNywiXFxtdV9BIl0sWzQsNywiMV97VEF9IiwxXV0=
\[\begin{tikzcd}
	{T^3 A} & {T^2A} && TA & {T^2A} \\
	{T^2A} & TA && {T^2 A} & TA
	\arrow["{\mu_{TA}}", from=1-1, to=1-2]
	\arrow["{T\mu_A}"', from=1-1, to=2-1]
	\arrow["{\mu_A}", from=1-2, to=2-2]
	\arrow["{\mu_A}"', from=2-1, to=2-2]
	\arrow["{\eta_{TA}}", from=1-4, to=1-5]
	\arrow["{T_{\eta A}}"', from=1-4, to=2-4]
	\arrow["{\mu_A}"', from=2-4, to=2-5]
	\arrow["{\mu_A}", from=1-5, to=2-5]
	\arrow["{1_{TA}}"{description}, from=1-4, to=2-5]
\end{tikzcd}\]
    $\eta$ is the unit, and $\mu$ is the multiplication of the monad.
\end{definition}
This definition, however, can easily be generalised using 2-categories.

\begin{definition}[Monad \cite{STREET1972149} \cite{Sobociński:2017}]
\label{def:monad-gen}
Let $\mathcal{C}$ denote any bicategory, strict or otherwise.
A \textbf{\underline{monad}} $(X,S)$ consists of
\begin{itemize}
  \item A 0-cell $X$
  \item A 1-cell $S : X \rightarrow X$
  \item A pair of 2-cells $\eta: 1_S \rightarrow S, \mu : SS \rightarrow S$
  called the \textit{unit} and \textit{multiplication} of the monad.
\end{itemize}
such that the following diagrams commute, up to a natural isomorphism.
% https://q.uiver.app/#q=WzAsOCxbMCwwLCJTXjMiXSxbMSwwLCJTXjIiXSxbMCwxLCJTXjIiXSxbMSwxLCJTIl0sWzMsMCwiUyJdLFs0LDAsIlNeMiJdLFszLDEsIlNeMiJdLFs0LDEsIlMiXSxbMCwxLCJcXG11IFMiXSxbMCwyLCJTXFxtdSIsMl0sWzEsMywiXFxtdSJdLFsyLDMsIlxcbXUiLDJdLFs0LDUsIlxcZXRhIFMiXSxbNCw2LCJTXFxldGEiLDJdLFs2LDcsIlxcbXUiLDJdLFs1LDcsIlxcbXUiXSxbNCw3LCIxX1MiLDFdXQ==
\[\begin{tikzcd}
	{S^3} & {S^2} && S & {S^2} \\
	{S^2} & S && {S^2} & S
	\arrow["{\mu S}", from=1-1, to=1-2]
	\arrow["S\mu"', from=1-1, to=2-1]
	\arrow["\mu", from=1-2, to=2-2]
	\arrow["\mu"', from=2-1, to=2-2]
	\arrow["{\eta S}", from=1-4, to=1-5]
	\arrow["S\eta"', from=1-4, to=2-4]
	\arrow["\mu"', from=2-4, to=2-5]
	\arrow["\mu", from=1-5, to=2-5]
	\arrow["{1_S}"{description}, from=1-4, to=2-5]
\end{tikzcd}\]
where the following notation is used.
\[S^2 := S\circ S := SS, \quad S^3 := S \circ S^2 := SSS\]
\[S\mu := S \triangleleft \mu, \quad \mu S := \mu \triangleright S\]
\end{definition}
\begin{definition}
  If we let the strict 2-category \textbf{Cat} of categories, functors, and natural transformations
  be the bicategory in Def. \ref{def:monad-gen}, then we get Def. \ref{def:monad}, where
  \begin{itemize}
    \item The 0-cell is the category $\mathcal{C}$ in Def. \ref{def:monad}
    \item The 1-cell is the endofunctor on $T: \mathcal{C} \rightarrow \mathcal{C}$
    \item The 2-cells are the natural transformations $\mu: T^2 \rightarrow T, \eta: 1_T \rightarrow T$.
  \end{itemize}
\end{definition}
\begin{definition}[Span \cite{nlab:span} \cite{Sobociński:2017}]
  A span in a category $\mathcal{C}$ from an object $A$ to $B$ is a diagram of the form
\[A \xleftarrow[]{f} C \xrightarrow[]{g} B\]
where $C$ is some other object in $\mathcal{C}$.

If $\mathcal{C}$ has pullbacks \cite{Abramsky_2010}, we can compose spans, 
so given a span from $A$ to $B$ and $B$ to $C$,
we can take a pullback and obtain a span from $A$ to $C$.
% https://q.uiver.app/#q=WzAsOSxbMiwwLCJEXFx0aW1lc19CIEUiXSxbMSwxLCJEIl0sWzMsMSwiRSJdLFswLDIsIkEiXSxbMiwyLCJCIl0sWzQsMiwiQyJdLFs3LDAsIkRcXHRpbWVzX0IgRSJdLFs2LDEsIkEiXSxbOCwxLCJDIl0sWzAsMSwicF9EIiwyXSxbMCwyLCJwX0UiXSxbMSwzLCJmIiwyXSxbMSw0LCJnIiwyXSxbMiw0LCJoIl0sWzIsNSwiayJdLFs2LDcsImZwX0QiLDJdLFs2LDgsImtwX0UiXV0=
\[\begin{tikzcd}
	&& {D\times_B E} &&&&& {D\times_B E} \\
	& D && E &&& A && C \\
	A && B && C
	\arrow["{p_D}"', from=1-3, to=2-2]
	\arrow["{p_E}", from=1-3, to=2-4]
	\arrow["f"', from=2-2, to=3-1]
	\arrow["g"', from=2-2, to=3-3]
	\arrow["h", from=2-4, to=3-3]
	\arrow["k", from=2-4, to=3-5]
	\arrow["{fp_D}"', from=1-8, to=2-7]
	\arrow["{kp_E}", from=1-8, to=2-9]
\end{tikzcd}\]
Maps between spans are then functions $f : C \rightarrow C'$ which make the diagram commute.
% https://q.uiver.app/#q=WzAsNCxbMSwxLCJDIl0sWzAsMiwiQSJdLFsyLDIsIkIiXSxbMSwwLCJDJyJdLFswLDEsImciXSxbMCwyLCJoIiwyXSxbMywxLCJnJyIsMix7ImN1cnZlIjoxfV0sWzMsMiwiZyciLDAseyJjdXJ2ZSI6LTF9XSxbMCwzLCJmIiwyXV0=
\[\begin{tikzcd}
	& {C'} \\
	& C \\
	A && B
	\arrow["g", from=2-2, to=3-1]
	\arrow["h"', from=2-2, to=3-3]
	\arrow["{g'}"', curve={height=6pt}, from=1-2, to=3-1]
	\arrow["{h'}", curve={height=-6pt}, from=1-2, to=3-3]
	\arrow["f"', from=2-2, to=1-2]
\end{tikzcd}\]
The bicategory Span$(\mathcal{C})$ then consists of
\begin{itemize}
  \item Objects of $\mathcal{C}$ as 0-cells
  \item Spans as 1-cells
  \item Maps between spans as 2-cells
\end{itemize}
\end{definition}
  Using the bicategory Span(\textbf{Set}) in Def. \ref{def:monad-gen} for a monad, we get the following.
\begin{example}[\cite{Sobociński:2017}]
\[\]
\begin{itemize}
  \item The 0-cell is a set $X$
  \item The endo-1-cell is a span $S$ of the form
  \[X\xleftarrow[]{\delta_0} Y \xrightarrow[]{\delta_1} X\]
  \item The unit 2-cell is the function 
  \[\eta: X\rightarrow Y, \quad \delta_0 \circ \eta (x) = \delta_1 \circ \eta (x) =x\]
  \item The multiplication 2-cell $SS \rightarrow S$ is the function
  \[\mu: Y\times_X Y \rightarrow Y, \quad \delta_0\circ \mu (y,y') = \delta_0(y), 
  \quad \delta_1\circ \mu (y,y') = \delta_1 (y')\]
  where $Y\times_X Y = \{(y,y') | \delta_1(y) = \delta_0(y')\}$.
\end{itemize}
which satisfy the axioms of
\begin{itemize}
  \item Associativity
  \[\mu(\mu(y,y'), y'') = \mu(y, \mu(y', y''))\]
  \item Unit
  \[\mu(y, \eta\circ \delta_1 (y)) = y = \mu(\eta \circ \delta_0 (y), y)\]
\end{itemize}
\end{example}
Here the monad of Span(\textbf{Set}) can be described as a category $\mathcal{C}$ where
\begin{itemize}
  \item $X$ is the set of objects
  \item $Y$ is the set of arrows
  \item $\delta_0,\delta_1 : Y\to X$ give the domain and codomain of each arrow respectively.
  \item The function $\eta:X\rightarrow Y$ gives the identity arrows $\eta(x) = 1_x$ for each $x\in X$.
  \item The hom-set for objects $a,b\in X$ is defined by 
  \[\mathcal{C}(a,b) := \{y\in Y | \delta_0 (y) = a \wedge \delta_1(y) = b \}\]
  \item $Y\times_X Y$ is the set of composable arrows, and $\mu: Y\times_X Y\to Y$ defines composition.
  \item The associative and unit axioms are defined in the monad as they should be in the category.
\end{itemize}
\begin{definition}[Locally small category \cite{Abramsky_2010} \cite{nlab:locally_small_category}]
  A \textbf{locally small category} is a category where all hom-sets are sets.
\end{definition}
\begin{definition}[Small category \cite{nlab:locally_small_category}]
  A \textbf{small category} is a locally small category where its set of objects is a set.
\end{definition}
\begin{example}
  The category of sets and functions, \textbf{Set}, is a locally small, but not small, category.
\end{example}
Then, we have that the monad on the bicategory of Span(\textbf{Set}) describes small categories,
since its set of objects is a set, and its hom-sets are sets.

\subsection{… , a monad is a monoid, …}
At the same time, the concept of a monoid on sets $(M,\cdot,1_M)$ can be generalised to categories.

% From Def. \ref{def:moncat}, 
% the monoid on $(\textbf{Set}, \times, \{*\})$
% gives rise to the "monoid-as-a-set" from Def. \ref{def:monoid}.

% Here we give definitions to aid in understanding the "monoid-as-a-set". 
\begin{note}
  Since a strict monoidal category $(\mathcal{C},\otimes, I, \alpha, \lambda, \rho)$ 
  (Def. \ref{def:strict-mon}) is such that $\alpha, \lambda, \rho$ are identities, 
  it can be described succinctly as $(\mathcal{C},\otimes, I)$.
\end{note}
\begin{definition}[Monoid \cite{Kissinger19:2024}]
  A \textbf{monoid} in a strict monoidal category $(\mathcal{C},\otimes, I)$ is an object $X$
  with morphisms 
  \[\mu : X\otimes X \rightarrow X, \quad \eta: I\rightarrow X\]
  satisfying the commuting diagrams describing associativity and unit axioms respectively.
% https://q.uiver.app/#q=WzAsOCxbMCwwLCJYXFxvdGltZXMgWFxcb3RpbWVzIFgiXSxbMSwwLCJYXFxvdGltZXMgWCJdLFsxLDEsIlgiXSxbMCwxLCJYXFxvdGltZXMgWCJdLFszLDAsIlgiXSxbNCwwLCJYXFxvdGltZXMgWCJdLFs0LDEsIlgiXSxbMywxLCJYXFxvdGltZXMgWCJdLFswLDEsIjFfWFxcb3RpbWVzIFxcbXUiXSxbMSwyLCJcXG11Il0sWzAsMywiXFxtdSBcXG90aW1lcyAxX1giLDJdLFszLDIsIlxcbXUiLDJdLFs0LDUsIlxcZXRhIFxcb3RpbWVzIDFfWCJdLFs1LDYsIlxcbXUiXSxbNCw2LCIxX1giLDFdLFs0LDcsIjFfWFxcb3RpbWVzIFxcZXRhIiwyXSxbNyw2LCJcXG11IiwyXV0=
\[\begin{tikzcd}
	{X\otimes X\otimes X} & {X\otimes X} && X & {X\otimes X} \\
	{X\otimes X} & X && {X\otimes X} & X
	\arrow["{1_X\otimes \mu}", from=1-1, to=1-2]
	\arrow["\mu", from=1-2, to=2-2]
	\arrow["{\mu \otimes 1_X}"', from=1-1, to=2-1]
	\arrow["\mu"', from=2-1, to=2-2]
	\arrow["{\eta \otimes 1_X}", from=1-4, to=1-5]
	\arrow["\mu", from=1-5, to=2-5]
	\arrow["{1_X}"{description}, from=1-4, to=2-5]
	\arrow["{1_X\otimes \eta}"', from=1-4, to=2-4]
	\arrow["\mu"', from=2-4, to=2-5]
\end{tikzcd}\]
\end{definition}

"Monoids-as-a-set" can then be formalised as follows.
\begin{example}["Monoids-as-a-set" \cite{Kissinger19:2024}]
Monoids in $(\textbf{Set}, \times, \{*\})$ are monoids 
% \begin{itemize}
%   \item Objects are a set $X$
%   \item Unit object is a one element set $\{*\}$
%   \item \[a\cdot b := \mu(a,b)\]
%   \item \[1 := \eta(*)\]
% \end{itemize}
\[(X, \mu : X\times X \rightarrow X, \eta : \{*\} \rightarrow X)\]
where
\[a\cdot b := \mu(a,b)\]
\[1 := \eta(*)\]
so $\mu$ is the binary relation, and $\eta(*)$ is the unit.
\end{example}
\begin{definition}[\cite{nlab:full_sub-2-category}]
  The 2-category $C$ is a full sub-2-category of $D$ if there is a 2-functor $F:C\to D$
  where for all objects $c_1,c_2 \in C$, the component functor $F_{c_1,c_2}$ is an equivalence of categories
  \[F_{c_1,c_2} : C(c_1,c_2) \xrightarrow[]{\simeq} D(F(c_1), F(c_2))\]
\end{definition}

Consider a strict 2-category $D$ and an object $c$ so that we have the category $D(c,c)$.
$D(c,c)$ is a full sub-2-category of $D$, since $D(c,c)$ is a 2-category with one object $c$.

\begin{proposition}
  $D(c,c)$ is a strict monoidal category.
\end{proposition}
\begin{proof}
\begin{itemize}
  \item $D(c,c)$ is the category, where objects are 1-cells, and arrows are 2-cells.
  \item The monoidal product $\otimes$ on 1-cells is composition of 1-cells.
  \item The unit object is the identity cell at $c, 1_c : c\to c$.
  \item 1-cells satisfy the associativity and unit axioms, by definition of the strict bicategory $D$,
  where the associator and unit natural isomorphisms are identities, so 
  \[(h \circ g) \circ f = h \circ (g \circ f), \quad 1 \circ f = f = f \circ 1\]
  \item 2-cells satisfy the associativity and unit axioms, by definition of arrows in the category $D(c,c)$
\end{itemize}
Thus, the natural transformation associator and unitors are identities, 
fulfilling the definition of a strict monoidal category.
\end{proof}
\begin{note}
  We denote $D(c,c)$ as $[C,C]$.
\end{note}
\begin{example}
  Then, a monoid on the strict monoidal category $([C,C], \circ, 1_C)$ is 
  an object $S : C \to C$ with morphisms
  \[\mu : S \circ S \to S, \quad \eta : 1_C \to S\]
  where the diagrams commute.
% https://q.uiver.app/#q=WzAsOCxbMCwwLCJTXjMiXSxbMSwwLCJTXjIiXSxbMCwxLCJTXjIiXSxbMSwxLCJTIl0sWzMsMCwiUyJdLFs0LDAsIlNeMiJdLFszLDEsIlNeMiJdLFs0LDEsIlMiXSxbMCwxLCJcXG11IFMiXSxbMCwyLCJTXFxtdSIsMl0sWzEsMywiXFxtdSJdLFsyLDMsIlxcbXUiLDJdLFs0LDUsIlxcZXRhIFMiXSxbNCw2LCJTXFxldGEiLDJdLFs2LDcsIlxcbXUiLDJdLFs1LDcsIlxcbXUiXSxbNCw3LCIxX1MiLDFdXQ==
\[\begin{tikzcd}
	{S^3} & {S^2} && S & {S^2} \\
	{S^2} & S && {S^2} & S
	\arrow["{\mu S}", from=1-1, to=1-2]
	\arrow["S\mu"', from=1-1, to=2-1]
	\arrow["\mu", from=1-2, to=2-2]
	\arrow["\mu"', from=2-1, to=2-2]
	\arrow["{\eta S}", from=1-4, to=1-5]
	\arrow["S\eta"', from=1-4, to=2-4]
	\arrow["\mu"', from=2-4, to=2-5]
	\arrow["\mu", from=1-5, to=2-5]
	\arrow["{1_S}"{description}, from=1-4, to=2-5]
\end{tikzcd}\]
This is precisely the definition of a monad in Def. \ref{def:monad-gen}, which is described with
\begin{itemize}
  \item The 0-cell $C$
  \item The 1-cell $S: C \to C$
  \item The 2-cells $\mu : S \circ S \to S, \quad \eta : 1_C \to S$
\end{itemize}
\end{example}
Then, a monad is a special kind of monoid on the category $([C,C], \circ, 1_C)$.

\begin{definition}[\cite{Kissinger20:2024}]
  A monad is a monoid in the category of endofunctors.
\end{definition}

\bibliographystyle{plain}
\bibliography{refs}

\end{document}
